% ============================================================
%  يُترجَم بـ XeLaTeX  (في Overleaf: Menu → Compiler → XeLaTeX)
%  الخط: Amiri (متوفر في Overleaf، أو حمّله من github.com/aliftype/amiri)
% ============================================================
\documentclass[12pt,a4paper]{article}
\usepackage[a4paper,margin=1.7cm,top=2.3cm,bottom=2.9cm,headheight=15pt,footskip=1.5cm]{geometry}
\usepackage{amsmath,amssymb}
\usepackage{icomma}
\usepackage[english,bidi=default,layout=lists]{babel}
\babelprovide[import,main]{arabic}
\babelfont{rm}[Scale=1.08]{Amiri}
\babelfont{sf}[Scale=1.08]{Amiri}
\usepackage[most]{tcolorbox}
\usepackage{tikz}
\usetikzlibrary{arrows.meta,calc,shadings}
\usepackage{array}
\usepackage{colortbl}
\usepackage{fancyhdr}
\usepackage{titlesec}
\usepackage{enumitem}
\usepackage{collcell}
\usepackage{needspace}
\usepackage{etoolbox}
\newcolumntype{C}{>{\collectcell\ar}c<{\endcollectcell}}
\pretocmd{\section}{\Needspace{15\baselineskip}}{}{}
\pretocmd{\subsection}{\Needspace{11\baselineskip}}{}{}

\definecolor{cDef}{HTML}{1F4E9C}
\definecolor{cThm}{HTML}{B3261E}
\definecolor{cRem}{HTML}{2E7D32}
\definecolor{cEx}{HTML}{C76B00}
\definecolor{cMeth}{HTML}{6A1B9A}
\definecolor{cSol}{HTML}{00695C}
\definecolor{cSec}{HTML}{1F3A68}

\titleformat{\section}{\Large\bfseries\color{cSec}}{\thesection.}{0.6em}{}[{\color{cSec!60}\titlerule[1pt]}]
\titleformat{\subsection}{\large\bfseries\color{cThm!85!black}}{\thesubsection.}{0.5em}{}
\titlespacing*{\section}{0pt}{14pt}{8pt}
\titlespacing*{\subsection}{0pt}{10pt}{5pt}
\setlength{\parindent}{0pt}
\setlength{\parskip}{4pt}
\setlist{itemsep=2pt,topsep=3pt}

\newtcolorbox{arbox}[3][]{enhanced,
  colback=#2!5,colframe=#2!80!black,arc=2mm,boxrule=0.8pt,
  left=6pt,right=6pt,top=8pt,bottom=5pt,
  fonttitle=\bfseries,coltitle=white,
  title={\foreignlanguage{arabic}{#3}},
  attach boxed title to top right={xshift=-5mm,yshift=-3mm},
  boxed title style={colback=#2!80!black,arc=1.5mm,boxrule=0pt},
  before upper={\selectlanguage{arabic}},
  before skip=8pt,after skip=8pt,breakable,#1}
\newenvironment{BOX}[2]{\begin{otherlanguage}{english}\begin{arbox}{#1}{#2}}{\end{arbox}\end{otherlanguage}}
\newenvironment{tarif}[1][تعريف]{\begin{BOX}{cDef}{#1}}{\end{BOX}}
\newenvironment{khassia}[1][خاصية]{\begin{BOX}{cThm}{#1}}{\end{BOX}}
\newenvironment{molahada}[1][ملاحظة]{\begin{BOX}{cRem}{#1}}{\end{BOX}}
\newenvironment{mithal}[1][مثال]{\begin{BOX}{cEx}{#1}}{\end{BOX}}
\newenvironment{tariqa}[1][طريقة]{\begin{BOX}{cMeth}{#1}}{\end{BOX}}

\newcounter{exo}
\newenvironment{tamrin}[1]{\stepcounter{exo}\begin{BOX}{cSec}{التمرين \theexo: #1}}{\end{BOX}}
\newenvironment{hal}{\Needspace{12\baselineskip}\begin{otherlanguage}{english}\begin{arbox}[breakable]{cSol}{حل التمرين \theexo}}{\end{arbox}\end{otherlanguage}}

\newenvironment{LTR}{\begin{otherlanguage}{english}}{\end{otherlanguage}}
\newcommand{\ar}[1]{\foreignlanguage{arabic}{#1}}

\renewcommand{\dots}{\mathinner{\ldotp\ldotp\ldotp}}
\renewcommand{\ldots}{\mathinner{\ldotp\ldotp\ldotp}}
\renewcommand{\cdots}{\mathinner{\cdotp\cdotp\cdotp}}
\newcommand{\N}{\mathbb{N}}
\newcommand{\Z}{\mathbb{Z}}
\newcommand{\D}{\mathbb{D}}
\newcommand{\Q}{\mathbb{Q}}
\newcommand{\R}{\mathbb{R}}
\newcommand{\pgcd}{\operatorname{PGCD}}
\newcommand{\ppcm}{\operatorname{PPCM}}

% \NLseg{a}{b}{[}{]} ; \NLright{a}{[} ; \NLleft{b}{]}  (interval on a number line)
\newcommand{\NLseg}[4]{\begin{LTR}\begin{tikzpicture}[baseline=-0.6ex]
 \draw[-{Stealth[length=2mm]},gray!70](0,0)--(3.4,0);
 \draw[line width=2pt,cThm!80!black](0.7,0)--(2.7,0);
 \node[cThm!80!black,font=\LARGE] at (0.7,0){$#3$};
 \node[cThm!80!black,font=\LARGE] at (2.7,0){$#4$};
 \node[below=3pt,font=\footnotesize] at (0.7,0){$#1$};
 \node[below=3pt,font=\footnotesize] at (2.7,0){$#2$};
\end{tikzpicture}\end{LTR}}
\newcommand{\NLright}[2]{\begin{LTR}\begin{tikzpicture}[baseline=-0.6ex]
 \draw[-{Stealth[length=2mm]},gray!70](0,0)--(3.4,0);
 \draw[line width=2pt,cThm!80!black](1.3,0)--(3.3,0);
 \node[cThm!80!black,font=\LARGE] at (1.3,0){$#2$};
 \node[below=3pt,font=\footnotesize] at (1.3,0){$#1$};
\end{tikzpicture}\end{LTR}}
\newcommand{\NLleft}[2]{\begin{LTR}\begin{tikzpicture}[baseline=-0.6ex]
 \draw[-{Stealth[length=2mm]},gray!70](0,0)--(3.4,0);
 \draw[line width=2pt,cThm!80!black](0.05,0)--(2.1,0);
 \node[cThm!80!black,font=\LARGE] at (2.1,0){$#2$};
 \node[below=3pt,font=\footnotesize] at (2.1,0){$#1$};
\end{tikzpicture}\end{LTR}}

% ---------- social-media icons (pure TikZ) ----------
\definecolor{cFB}{HTML}{1877F2}
\definecolor{cTG}{HTML}{229ED9}
\definecolor{cIG1}{HTML}{FEDA75}
\definecolor{cIG2}{HTML}{D62976}
\definecolor{cIG3}{HTML}{4F5BD5}
\newcommand{\fbicon}{\begin{LTR}\begin{tikzpicture}[baseline=-0.75ex,x=1em,y=1em]
 \fill[cFB] (0,0) circle (0.55);
 \node[white,inner sep=0,font=\usefont{TU}{lmss}{bx}{n}\small] at (0.04,-0.02) {f};
\end{tikzpicture}\end{LTR}}
\newcommand{\igicon}{\begin{LTR}\begin{tikzpicture}[baseline=-0.75ex,x=1em,y=1em]
 \shade[lower left=cIG1,lower right=cIG2,upper left=cIG2,upper right=cIG3,rounded corners=0.28em] (-0.55,-0.55) rectangle (0.55,0.55);
 \draw[white,line width=0.09em,rounded corners=0.18em] (-0.36,-0.36) rectangle (0.36,0.36);
 \draw[white,line width=0.09em] (0,0) circle (0.17);
 \fill[white] (0.22,0.22) circle (0.05);
\end{tikzpicture}\end{LTR}}
\newcommand{\tgicon}{\begin{LTR}\begin{tikzpicture}[baseline=-0.75ex,x=1em,y=1em]
 \fill[cTG] (0,0) circle (0.55);
 \fill[white] (-0.33,0.0) -- (0.32,0.27) -- (0.2,-0.3) -- (0.0,-0.15) -- (-0.08,-0.28) -- (-0.08,-0.08) -- cycle;
 \fill[cTG!70!black] (-0.08,-0.08) -- (0.25,0.2) -- (0.0,-0.15) -- (-0.08,-0.28) -- cycle;
\end{tikzpicture}\end{LTR}}
% page number: black box, white two-digit number, at the right of the social row
\newcommand{\PageBox}{\begin{LTR}\tikz[baseline=(p.base)]\node[fill=black,text=white,inner sep=2.5pt,minimum width=1.6em,minimum height=1.3em,font=\usefont{TU}{lmss}{bx}{n}\small](p){\ifnum\value{page}<10 0\fi\arabic{page}};\end{LTR}}
% visual left-to-right order: Telegram | Instagram | Facebook | page box
\newcommand{\SocialRow}{\begin{otherlanguage}{english}%
 \ar{\textbf{الأستاذ سوايسية هشام للرياضيات}}\ \tgicon\hfill
 \textbf{prof\_hicham\_math}\ \igicon\hfill
 \ar{\textbf{الأستاذ سوايسية هشام للرياضيات}}\ \fbicon\hspace{2.2em}\PageBox\end{otherlanguage}}

\pagestyle{fancy}
\fancyhf{}
\renewcommand{\headrulewidth}{0.6pt}
\renewcommand{\footrulewidth}{0.4pt}
% \SetHeader{right head: document title}{left head: level -- stream}
\newcommand{\SetHeader}[2]{%
 \fancyhead[C]{\parbox{\textwidth}{\ar{#1\hfill #2}}}%
 \fancyfoot[C]{\parbox{\textwidth}{\small\SocialRow}}}
\usepackage{tkz-tab}
\usepackage{multicol}
% variation / sign table (LTR)
\newenvironment{TAB}{\par\begin{center}\begin{LTR}\begin{tikzpicture}[scale=0.95,every node/.style={font=\small}]}{\end{tikzpicture}\end{LTR}\end{center}}
% grid + axes, clipped: \begin{GR}{xmin}{xmax}{ymin}{ymax}{unit cm}
\newenvironment{GR}[5]{\par\begin{center}\begin{LTR}\begin{tikzpicture}[x=#5cm,y=#5cm,font=\footnotesize]
 \draw[gray!25,very thin] (#1,#3) grid (#2,#4);
 \draw[-{Stealth[length=2mm]}] (#1,0)--(#2+0.3,0) node[below left]{$x$};
 \draw[-{Stealth[length=2mm]}] (0,#3)--(0,#4+0.3) node[below left]{$y$};
 \node[below left] at (0,0){$O$};
 \begin{scope}\clip (#1,#3) rectangle (#2,#4);}
{\end{scope}\end{tikzpicture}\end{LTR}\end{center}}
\newcommand{\Ticks}[2]{\foreach \t in {#1} \node[below,font=\scriptsize] at (\t,0){$\t$}; \foreach \t in {#2} \node[left,font=\scriptsize] at (0,\t){$\t$};}
\newcommand{\bx}[1]{\ensuremath{\colorbox{cDef!12}{$\displaystyle #1$}}}
\newcommand{\TwoCol}[4]{\par\begin{LTR}\noindent\begin{minipage}[t]{#1\textwidth}\selectlanguage{arabic}#2\end{minipage}\hfill\begin{minipage}[t]{#3\textwidth}\selectlanguage{arabic}#4\end{minipage}\end{LTR}\par}
\newcommand{\RL}[4]{\TwoCol{#3}{#4}{#1}{#2}}  % \RL{w1}{right content}{w2}{left content}
\newcommand{\ok}{$\textcolor{cRem}{\checkmark}$}
\newcommand{\limn}{\displaystyle\lim_{n\to+\infty}}
% sub-questions, each on its own line
\newenvironment{sq}{\begin{itemize}[leftmargin=2.4em,itemsep=3pt]}{\end{itemize}}
\newcommand{\qa}{\item[(أ)]}\newcommand{\qb}{\item[(ب)]}\newcommand{\qc}{\item[(جـ)]}\newcommand{\qd}{\item[(د)]}\newcommand{\qe}{\item[(هـ)]}
\newcommand{\pa}{\par\smallskip\textbf{(أ)} }\newcommand{\pb}{\par\smallskip\textbf{(ب)} }\newcommand{\pc}{\par\smallskip\textbf{(جـ)} }\newcommand{\pd}{\par\smallskip\textbf{(د)} }\newcommand{\pe}{\par\smallskip\textbf{(هـ)} }
\newcommand{\QQ}[1]{\par\smallskip\textbf{#1.} }
\newenvironment{tmr}{\stepcounter{exo}\begin{BOX}{cSec}{تمرين \theexo}}{\end{BOX}}
\newcommand{\ABC}[3]{\\[2pt]\begin{tabular}{p{0.28\textwidth}p{0.28\textwidth}p{0.28\textwidth}}(جـ) \ #3 & (ب) \ #2 & (أ) \ #1\end{tabular}}
\pretocmd{\section}{\Needspace{24\baselineskip}}{}{}
\newtcbox{\bxtbox}{on line,colback=cDef!12,colframe=cDef!12,boxrule=0pt,arc=1pt,boxsep=0pt,left=2pt,right=2pt,top=1pt,bottom=1pt}
\newcommand{\bxt}[1]{\foreignlanguage{english}{\bxtbox{\ar{#1}}}}
\SetHeader{\small\bfseries\color{cSec}تمارين تجريبية: الدالة اللوغاريتمية}{\small\color{cSec}السنة الثالثة ثانوي -- الشعب العلمية}
\newenvironment{FIG}{\par\begin{center}\begin{LTR}}{\end{LTR}\end{center}}
\begin{document}
\begin{LTR}
\begin{tcolorbox}[enhanced,colback=cSec,colframe=cSec,arc=3mm,halign=center,boxrule=0pt,top=12pt,bottom=12pt,before upper={\selectlanguage{arabic}}]
{\color{white}\fontsize{26}{32}\selectfont\bfseries تمارين تجريبية في الدالة اللوغاريتمية}\\[8pt]
{\color{white}\large تمارين متدرجة -- بدون حلول (للتدرّب)}\\[5pt]
{\color{white}\large السنة الثالثة ثانوي -- الشعب العلمية}
\end{tcolorbox}
\end{LTR}
\section{مجموعة التعريف والتبسيط}

% ------------------------------------------------------------------
\begin{tamrin}{مجموعة تعريف دالة}
عيّن مجموعة تعريف كل دالة من الدوال التالية:
\QQ{1} $f_1(x)=\ln(2x-4)$
\QQ{2} $f_2(x)=\ln(3-x)+\ln(x+1)$
\QQ{3} $f_3(x)=\ln\left(\dfrac{x-1}{x+2}\right)$
\QQ{4} $f_4(x)=\ln\left|x^2-1\right|$
\QQ{5} $f_5(x)=\dfrac{x}{\ln x}$
\QQ{6} $f_6(x)=\dfrac{1}{1-\ln x}$
\QQ{7} $f_7(x)=\ln\left(e^{x}-1\right)$
\end{tamrin}

% ------------------------------------------------------------------
\begin{tamrin}{تبسيط عبارات}
\QQ{1} اكتب بدلالة $\ln 2$ و $\ln 3$ كلا من الأعداد:
\begin{sq}
\qa $A=\ln 8+\ln 2-\ln 4$
\qb $B=\ln 72$
\qc $C=\ln\sqrt{12}$
\qd $D=\ln\dfrac{4}{9}$
\qe $E=\ln 27-2\ln 3+\ln\dfrac{1}{9}$
\end{sq}
\QQ{2} بسّط ما يلي:
\begin{sq}
\qa $F=\ln\sqrt{e}+\ln\dfrac{1}{e^{2}}-\ln\left(e^{3}\right)$
\qb $G=\ln\left(\sqrt{3}+\sqrt{2}\right)+\ln\left(\sqrt{3}-\sqrt{2}\right)$
\qc $H=e^{2\ln 3}-\ln\left(e^{-2}\right)$
\end{sq}
\QQ{3} بيّن أنه من أجل كل عدد حقيقي $x$: \ $\ln\left(e^{x}+1\right)-x=\ln\left(1+e^{-x}\right)$.
\end{tamrin}

% ------------------------------------------------------------------
\begin{tamrin}{خواص جبرية وتطبيقات}
\QQ{1} بيّن أنه من أجل كل عدد حقيقي $x$ من $\left]1;+\infty\right[$:
\ $\ln\left(x^{2}-1\right)-2\ln x=\ln\left(1-\dfrac{1}{x^{2}}\right)$.
\QQ{2} نعتبر الدالة $g$ المعرفة بـ: $g(x)=\ln\left(x+\sqrt{x^{2}+1}\right)$.
\begin{sq}
\qa بيّن أنه من أجل كل عدد حقيقي $x$: $x+\sqrt{x^{2}+1}>0$، ثم استنتج مجموعة تعريف $g$.
\qb احسب $\left(x+\sqrt{x^{2}+1}\right)\left(-x+\sqrt{x^{2}+1}\right)$.
\qc استنتج أن الدالة $g$ فردية.
\end{sq}
\QQ{3} احسب المجموع: $S=\ln\dfrac12+\ln\dfrac23+\ln\dfrac34+\cdots+\ln\dfrac{99}{100}$، ثم اكتبه بدلالة $\ln 10$.
\end{tamrin}

\section{المعادلات والمتراجحات والجمل}

% ------------------------------------------------------------------
\begin{tamrin}{معادلات تتضمن اللوغاريتم}
حل في $\R$ المعادلات التالية:
\QQ{1} $\ln(2x-1)=0$
\QQ{2} $\ln(x+3)=\ln(5-x)$
\QQ{3} $\ln x+\ln(x-2)=\ln 8$
\QQ{4} $2\ln x=\ln(x+6)$
\QQ{5} $\ln|x-1|=2$
\end{tamrin}

% ------------------------------------------------------------------
\begin{tamrin}{متراجحات وجداول الإشارة}
\QQ{1} حل في $\R$ المتراجحات التالية:
\begin{sq}
\qa $\ln(3x-2)\leq 0$
\qb $\ln(x+1)\geq\ln(2x-3)$
\qc $\ln x+\ln(x+2)<\ln 3$
\end{sq}
\QQ{2} نعتبر الدالة $g$ المعرفة على $\left]0;+\infty\right[$ بـ: $g(x)=(x-2)\ln x$.
\begin{sq}
\qa ادرس إشارة $g(x)$ في جدول.
\qb استنتج حلول المتراجحة $g(x)\leq 0$.
\end{sq}
\QQ{3} نعتبر الدالة $h$ المعرفة بـ: $h(x)=\dfrac{\ln x}{1-\ln x}$.
\begin{sq}
\qa عيّن مجموعة تعريف $h$.
\qb ادرس إشارة $h(x)$، ثم حل المتراجحة $h(x)>0$.
\end{sq}
\end{tamrin}

% ------------------------------------------------------------------
\begin{tamrin}{تغيير المجهول $t=\ln x$}
حل في $\left]0;+\infty\right[$ ما يلي:
\QQ{1} $\left(\ln x\right)^{2}-3\ln x+2=0$
\QQ{2} $2\left(\ln x\right)^{2}+\ln x-1=0$
\QQ{3} $\left(\ln x\right)^{3}-4\ln x=0$
\QQ{4} $\left(\ln x\right)^{2}-2\ln x-3\leq 0$
\QQ{5} $2\left(\ln x\right)^{2}-5\ln x+2>0$
\end{tamrin}

% ------------------------------------------------------------------
\begin{tamrin}{جمل معادلات تتضمن اللوغاريتم}
حل في $\R^{2}$ الجمل التالية:
\QQ{1} $\begin{cases} x+y=7\\ \ln x+\ln y=\ln 12\end{cases}$

\QQ{2} $\begin{cases} \ln x+2\ln y=4\\ 3\ln x-\ln y=5\end{cases}$
\QQ{3} $\begin{cases} x^{2}+y^{2}=20\\ \ln x+\ln y=3\ln 2\end{cases}$

\QQ{4} $\begin{cases} e^{x}\times e^{y}=e^{5}\\ \ln x+\ln y=\ln 6\end{cases}$
\end{tamrin}

% ------------------------------------------------------------------
\begin{tamrin}{معادلات ومتراجحات تمزج بين $e^{x}$ و $\ln$}
حل في $\R$ ما يلي:
\QQ{1} $e^{2x+1}>3$
\QQ{2} $e^{2x}-4e^{x}+3=0$
\QQ{3} $e^{x}-6e^{-x}-1=0$
\QQ{4} $e^{2x}-5e^{x}+6\leq 0$
\QQ{5} $\ln\left(e^{x}-1\right)=x-\ln 2$
\end{tamrin}

\section{النهايات}

% ------------------------------------------------------------------
\begin{tamrin}{نهايات مباشرة}
احسب النهايات التالية:
\QQ{1} $\displaystyle\lim_{x\to+\infty}\ln(2x+1)$
\QQ{2} $\displaystyle\lim_{x\to 1^{+}}\ln(x-1)$
\QQ{3} $\displaystyle\lim_{x\to 0^{+}}\left(x-\ln x\right)$
\QQ{4} $\displaystyle\lim_{x\to 0^{+}}\frac{1}{\ln x}$
\QQ{5} $\displaystyle\lim_{x\to+\infty}\ln\left(\frac{2x+1}{x+3}\right)$
\QQ{6} $\displaystyle\lim_{x\to 2^{+}}\ln\left(x^{2}-4\right)$
\end{tamrin}

% ------------------------------------------------------------------
\begin{tamrin}{النهايات الشهيرة}
\textbf{نذكّر بالنهايات:}
$\displaystyle\lim_{x\to+\infty}\frac{\ln x}{x}=0$،\ \
$\displaystyle\lim_{x\to 0^{+}}x\ln x=0$،\ \
$\displaystyle\lim_{x\to 0}\frac{\ln(1+x)}{x}=1$،\ \
$\displaystyle\lim_{x\to 1}\frac{\ln x}{x-1}=1$.

باستعمالها احسب ما يلي:
\QQ{1} $\displaystyle\lim_{x\to+\infty}\frac{\ln x}{x^{2}}$
\QQ{2} $\displaystyle\lim_{x\to 0^{+}}x^{2}\ln x$
\QQ{3} $\displaystyle\lim_{x\to 0}\frac{\ln(1+3x)}{x}$
\QQ{4} $\displaystyle\lim_{x\to 1}\frac{\ln x}{x^{2}-1}$
\QQ{5} $\displaystyle\lim_{x\to 0^{+}}\sqrt{x}\,\ln x$
\QQ{6} $\displaystyle\lim_{x\to e}\frac{\ln x-1}{x-e}$
\end{tamrin}

% ------------------------------------------------------------------
\begin{tamrin}{إزالة حالات عدم التعيين}
احسب النهايات التالية بعد تحديد نوع حالة عدم التعيين:
\QQ{1} $\displaystyle\lim_{x\to+\infty}\left(x-\ln x\right)$
\QQ{2} $\displaystyle\lim_{x\to+\infty}\frac{\ln(x+1)}{x}$
\QQ{3} $\displaystyle\lim_{x\to 0^{+}}\left(\frac1x+\ln x\right)$
\QQ{4} $\displaystyle\lim_{x\to+\infty}\frac{\left(\ln x\right)^{2}}{x}$
\QQ{5} $\displaystyle\lim_{x\to+\infty}\left[\ln\left(x^{2}+1\right)-2\ln x\right]$
\QQ{6} $\displaystyle\lim_{x\to+\infty}\left(x\ln x-x\right)$
\end{tamrin}

% ------------------------------------------------------------------
\begin{tamrin}{تغيير المتغير والتفسير الهندسي}
\textbf{الجزء الأول:} باستعمال تغيير مناسب للمتغير، احسب:
\QQ{1} $\displaystyle\lim_{x\to+\infty}x\ln\left(1+\frac1x\right)$ \quad (ضع $t=\dfrac1x$)
\QQ{2} $\displaystyle\lim_{x\to 0^{+}}x\left(\ln x\right)^{2}$ \quad (ضع $t=\sqrt x$)
\QQ{3} $\displaystyle\lim_{x\to 1}\frac{\ln x}{\sqrt x-1}$ \quad (ضع $t=\sqrt x$)

\textbf{الجزء الثاني:} نعتبر الدالة $h$ المعرفة على $\left]0;+\infty\right[$ بـ: $h(x)=x-1+\dfrac{\ln x}{x}$، و $(C_h)$ تمثيلها البياني في معلم متعامد ومتجانس.
\QQ{1} احسب $\displaystyle\lim_{x\to 0^{+}}h(x)$ وفسّر النتيجة هندسيا.
\QQ{2} احسب $\displaystyle\lim_{x\to+\infty}h(x)$.
\QQ{3} بيّن أن المستقيم $(\Delta)$ ذا المعادلة $y=x-1$ مستقيم مقارب مائل لـ $(C_h)$ بجوار $+\infty$.
\QQ{4} ادرس الوضع النسبي لـ $(C_h)$ و $(\Delta)$.
\end{tamrin}

\section{الاشتقاق واتجاه التغير والمماس}

\begin{tamrin}{حساب المشتقات}
في كل حالة من الحالات التالية، الدالة $f$ قابلة للاشتقاق على المجال $I$. احسب $f'(x)$ من أجل كل عدد حقيقي $x$ من $I$:
\begin{sq}
\qa $f(x)=\ln\left(x^2+1\right)$ \quad و \quad $I=\R$.
\qb $f(x)=\ln(2x-3)$ \quad و \quad $I=\left]\frac32;+\infty\right[$.
\qc $f(x)=x\ln x-x$ \quad و \quad $I=]0;+\infty[$.
\qd $f(x)=\dfrac{\ln x}{x}$ \quad و \quad $I=]0;+\infty[$.
\qe $f(x)=(\ln x)^2$ \quad و \quad $I=]0;+\infty[$.
\item[(و)] $f(x)=\ln\left(\dfrac{x-1}{x+1}\right)$ \quad و \quad $I=]1;+\infty[$.
\item[(ز)] $f(x)=\ln\left(e^x+1\right)$ \quad و \quad $I=\R$.
\item[(ح)] $f(x)=\ln|x-2|$ \quad و \quad $I=]-\infty;2[$.
\end{sq}
\end{tamrin}

\begin{tamrin}{اتجاه التغير ومقارنة عددين}
نعتبر الدالة $f$ المعرفة على المجال $]0;+\infty[$ بـ: $f(x)=\dfrac{\ln x}{x}$، و $(C_f)$ تمثيلها البياني في معلم متعامد ومتجانس.
\QQ{1} احسب $\displaystyle\lim_{x\to0^{+}}f(x)$ و $\displaystyle\lim_{x\to+\infty}f(x)$، ثم فسّر النتيجتين بيانيًا.
\QQ{2} بيّن أنه من أجل كل عدد حقيقي $x$ من $]0;+\infty[$: $f'(x)=\dfrac{1-\ln x}{x^2}$، ثم شكّل جدول تغيرات الدالة $f$.
\QQ{3} استنتج أنه من أجل كل عدد حقيقي $x$ من $]0;+\infty[$: $\ln x\leqslant \dfrac{x}{e}$.
\QQ{4} عيّن نقطة تقاطع $(C_f)$ مع حامل محور الفواصل، ثم اكتب معادلة للمماس $(T)$ عندها.
\QQ{5} باستعمال اتجاه تغير $f$ على المجال $[e;+\infty[$، قارن بين العددين $e^{\pi}$ و $\pi^{e}$.
\end{tamrin}

\begin{tamrin}{معادلة المماس ومحور التناظر}
نعتبر الدالة $f$ المعرفة بـ: $f(x)=\ln\left(x^2-x+1\right)$، و $(C_f)$ تمثيلها البياني في معلم متعامد ومتجانس.
\QQ{1} بيّن أن $f$ معرفة على $\R$.
\QQ{2} بيّن أنه من أجل كل عدد حقيقي $x$: $f(1-x)=f(x)$. ماذا تستنتج بالنسبة إلى $(C_f)$؟
\QQ{3} احسب $f'(x)$ ثم ادرس اتجاه تغير الدالة $f$ وشكّل جدول تغيراتها (مع النهايتين).
\QQ{4} اكتب معادلة لكل من المماسين $(T_0)$ و $(T_1)$ لـ $(C_f)$ في النقطتين ذواتي الفاصلتين $0$ و $1$ على الترتيب.
\QQ{5} هل يوجد مماس لـ $(C_f)$ يوازي المستقيم $(\Delta):\ y=x$ غير $(T_1)$؟ إن وجد فاكتب معادلته.
\end{tamrin}

\begin{tamrin}{إثبات متباينات باستعمال دالة مساعدة}
\QQ{1} نعتبر الدالة $g$ المعرفة على $]0;+\infty[$ بـ: $g(x)=\ln x-x+1$.
\begin{sq}
\qa ادرس اتجاه تغير الدالة $g$ (دون حساب النهايات).
\qb استنتج أنه من أجل كل عدد حقيقي $x$ من $]0;+\infty[$: $\ln x\leqslant x-1$.
\qc فسّر النتيجة بيانيًا بالنسبة إلى منحنى الدالة $\ln$ ومماسه في النقطة ذات الفاصلة $1$.
\end{sq}
\QQ{2} بتعويض $x$ بـ $\dfrac1x$، بيّن أنه من أجل كل عدد حقيقي $x$ من $]0;+\infty[$: $1-\dfrac1x\leqslant\ln x$.
\QQ{3} استنتج أنه من أجل كل عدد حقيقي $x$ من $]0;+\infty[$: $\dfrac{x}{x+1}\leqslant\ln(1+x)\leqslant x$.
\QQ{4}
\begin{sq}
\qa بيّن أنه من أجل كل عدد حقيقي $x$ من $]0;+\infty[$: $\dfrac{x}{x+1}\leqslant x\ln\left(1+\dfrac1x\right)\leqslant 1$.
\qb استنتج $\displaystyle\lim_{x\to+\infty}x\ln\left(1+\frac1x\right)$ ثم $\displaystyle\lim_{x\to+\infty}\left(1+\frac1x\right)^{x}$.
\end{sq}
\QQ{5} نعتبر الدالة $h$ المعرفة على $[1;+\infty[$ بـ: $h(x)=\ln x-\dfrac{2(x-1)}{x+1}$.

بيّن أن $h'(x)=\dfrac{(x-1)^2}{x(x+1)^2}$، ثم استنتج أنه من أجل كل عدد حقيقي $x\geqslant1$: $\ln x\geqslant\dfrac{2(x-1)}{x+1}$.
\end{tamrin}

\section{دراسة دوال لوغاريتمية}

\begin{tamrin}{دالة مساعدة ومستقيم مقارب مائل ومناقشة بيانية}
\textbf{الجزء الأول:} نعتبر الدالة $g$ المعرفة على $]0;+\infty[$ بـ: $g(x)=x^2-1+\ln x$.
\QQ{1} ادرس اتجاه تغير الدالة $g$ ثم احسب $g(1)$.
\QQ{2} استنتج إشارة $g(x)$ على المجال $]0;+\infty[$.

\textbf{الجزء الثاني:} نعتبر الدالة $f$ المعرفة على $]0;+\infty[$ بـ: $f(x)=x+1-\dfrac{\ln x}{x}$، و $(C_f)$ تمثيلها البياني في معلم متعامد ومتجانس $(O;\vec i,\vec j)$.
\QQ{1} احسب $\displaystyle\lim_{x\to0^{+}}f(x)$ وفسّر النتيجة بيانيًا، ثم احسب $\displaystyle\lim_{x\to+\infty}f(x)$.
\QQ{2}
\begin{sq}
\qa بيّن أنه من أجل كل عدد حقيقي $x$ من $]0;+\infty[$: $f'(x)=\dfrac{g(x)}{x^2}$.
\qb استنتج اتجاه تغير الدالة $f$ ثم شكّل جدول تغيراتها.
\end{sq}
\QQ{3}
\begin{sq}
\qa بيّن أن المستقيم $(\Delta)$ ذا المعادلة $y=x+1$ مقارب مائل لـ $(C_f)$ بجوار $+\infty$.
\qb ادرس الوضع النسبي لـ $(C_f)$ و $(\Delta)$.
\end{sq}
\QQ{4} بيّن أن $(C_f)$ يقبل مماسًا $(T)$ يوازي $(\Delta)$ يطلب تعيين معادلة له.
\QQ{5} أنشئ $(\Delta)$ و $(T)$ و $(C_f)$.
\QQ{6} ناقش بيانيًا، حسب قيم الوسيط الحقيقي $m$، عدد حلول المعادلة $f(x)=m$.
\end{tamrin}

\begin{tamrin}{العدد $\alpha$ والحصر وتقاطع المنحنى مع محور الفواصل}
\textbf{الجزء الأول:} نعتبر الدالة $g$ المعرفة على $]0;+\infty[$ بـ: $g(x)=x^2+1+\ln x$.
\QQ{1} احسب نهايتي $g$ عند $0$ وعند $+\infty$، ثم ادرس اتجاه تغيرها وشكّل جدول تغيراتها.
\QQ{2} بيّن أن المعادلة $g(x)=0$ تقبل حلًا وحيدًا $\alpha$ حيث $0{,}32<\alpha<0{,}34$.
\QQ{3} استنتج إشارة $g(x)$ على $]0;+\infty[$.

\textbf{الجزء الثاني:} نعتبر الدالة $f$ المعرفة على $]0;+\infty[$ بـ: $f(x)=x-\dfrac{2+\ln x}{x}$، و $(C_f)$ تمثيلها البياني في معلم متعامد ومتجانس.
\QQ{1} احسب $\displaystyle\lim_{x\to0^{+}}f(x)$ و $\displaystyle\lim_{x\to+\infty}f(x)$.
\QQ{2} بيّن أنه من أجل كل عدد حقيقي $x$ من $]0;+\infty[$: $f'(x)=\dfrac{g(x)}{x^2}$، ثم شكّل جدول تغيرات $f$.
\QQ{3} بيّن أن $f(\alpha)=2\alpha-\dfrac1\alpha$ ثم استنتج حصرًا للعدد $f(\alpha)$.
\QQ{4}
\begin{sq}
\qa بيّن أن المستقيم $(\Delta):\ y=x$ مقارب مائل لـ $(C_f)$ بجوار $+\infty$.
\qb ادرس الوضع النسبي لـ $(C_f)$ و $(\Delta)$.
\end{sq}
\QQ{5} بيّن أن $(C_f)$ يقطع حامل محور الفواصل في نقطتين فاصلتاهما $x_1$ و $x_2$ حيث: $0{,}1<x_1<0{,}2$ و $1{,}5<x_2<1{,}6$.
\QQ{6} أنشئ $(\Delta)$ و $(C_f)$.
\end{tamrin}

\begin{tamrin}{مقارنة مع منحنى الدالة ln ومناقشة بيانية}
\textbf{الجزء الأول:} نعتبر الدالة $g$ المعرفة على $]0;+\infty[$ بـ: $g(x)=x+1+\ln x$.
\QQ{1} ادرس اتجاه تغير الدالة $g$.
\QQ{2} بيّن أن المعادلة $g(x)=0$ تقبل حلًا وحيدًا $\alpha$ حيث $0{,}27<\alpha<0{,}28$، ثم استنتج إشارة $g(x)$.

\textbf{الجزء الثاني:} نعتبر الدالة $f$ المعرفة على $]0;+\infty[$ بـ: $f(x)=\dfrac{x\ln x}{x+1}$، و $(C_f)$ تمثيلها البياني، و $(\Gamma)$ منحنى الدالة $x\mapsto\ln x$ في معلم متعامد ومتجانس.
\QQ{1} احسب $\displaystyle\lim_{x\to0^{+}}f(x)$ و $\displaystyle\lim_{x\to+\infty}f(x)$.
\QQ{2} بيّن أنه من أجل كل عدد حقيقي $x$ من $]0;+\infty[$: $f'(x)=\dfrac{g(x)}{(x+1)^2}$، ثم شكّل جدول تغيرات $f$.
\QQ{3} بيّن أن $f(\alpha)=-\alpha$، ثم استنتج حصرًا للعدد $f(\alpha)$.
\QQ{4}
\begin{sq}
\qa احسب $\displaystyle\lim_{x\to+\infty}\left[f(x)-\ln x\right]$ وفسّر النتيجة بيانيًا.
\qb ادرس الوضع النسبي لـ $(C_f)$ و $(\Gamma)$.
\end{sq}
\QQ{5} اكتب معادلة للمماس $(T)$ لـ $(C_f)$ في النقطة ذات الفاصلة $1$.
\QQ{6} أنشئ $(T)$ و $(\Gamma)$ و $(C_f)$.
\QQ{7} ناقش بيانيًا، حسب قيم الوسيط الحقيقي $m$، عدد حلول المعادلة $x\ln x=m(x+1)$ ذات المجهول $x$ من $]0;+\infty[$.
\end{tamrin}

\section{اللوغاريتم العشري ومسائل من الواقع}

\begin{tamrin}{خواص اللوغاريتم العشري}
نذكّر أن اللوغاريتم العشري هو الدالة $\log$ المعرفة على $]0;+\infty[$ بـ: $\log x=\dfrac{\ln x}{\ln10}$.
\QQ{1} احسب: $\log1000$ \quad و \quad $\log0{,}01$ \quad و \quad $\log\left(10^5\sqrt{10}\right)$.
\QQ{2} بسّط العددين: $A=\log50+\log2$ \quad و \quad $B=\log0{,}2+\log5-\log1$.
\QQ{3} حل في $\R$ المعادلة: $\log x+\log(x-3)=1$.
\QQ{4} حل في $\R$ المتراجحة: $\log(2x-1)\leqslant2$.
\QQ{5}
\begin{sq}
\qa بيّن أن عدد أرقام كتابة عدد طبيعي $N\geqslant1$ في النظام العشري هو $k$ حيث $10^{k-1}\leqslant N<10^{k}$.
\qb علمًا أن $\log2\approx0{,}30103$، عيّن عدد أرقام العدد $2^{2026}$.
\end{sq}
\QQ{6} عيّن أصغر عدد طبيعي $n$ حيث: $(0{,}9)^n<0{,}01$.
\end{tamrin}

\begin{tamrin}{حموضة محلول: الـ pH}
يُعرّف الـ pH لمحلول مائي بالعلاقة: $\text{pH}=-\log\left[\text{H}_3\text{O}^+\right]$، حيث $\left[\text{H}_3\text{O}^+\right]$ هو تركيز شوارد الهيدرونيوم بـ $\text{mol/L}$.

يكون المحلول حمضيًا إذا كان $\text{pH}<7$، معتدلًا إذا كان $\text{pH}=7$، وأساسيًا إذا كان $\text{pH}>7$.
\QQ{1} احسب الـ pH لمحلول تركيزه $\left[\text{H}_3\text{O}^+\right]=2{,}5\times10^{-4}\ \text{mol/L}$ (تدوّر النتيجة إلى $10^{-2}$). ما طبيعة هذا المحلول؟
\QQ{2} محلول له $\text{pH}=8{,}2$. احسب تركيزه $\left[\text{H}_3\text{O}^+\right]$، وما طبيعته؟
\QQ{3} نعتبر الدالة $p$ المعرفة على $]0;+\infty[$ بـ: $p(x)=-\log x$. ادرس اتجاه تغيرها وفسّر النتيجة كيميائيًا.
\QQ{4} نمدّد محلولًا حمضيًا فيُقسم تركيزه على $10$. بيّن أن الـ pH يزداد بـ $1$.
\QQ{5} نضرب تركيز محلول في عدد $k>0$. بيّن أن الـ pH يتغير بـ $-\log k$، ثم عيّن قيمة مقربة لـ $k$ حتى ينقص الـ pH بـ $0{,}3$.
\end{tamrin}

\begin{tamrin}{شدة الصوت والنمو السكاني}
\textbf{الجزء الأول (الديسيبل):} يُعطى مستوى الشدة الصوتية $L$ (بالديسيبل $\text{dB}$) لصوت شدته $I$ (بـ $\text{W/m}^2$) بالعلاقة:
\[L=10\log\left(\frac{I}{I_0}\right)\qquad\text{حيث}\qquad I_0=10^{-12}\ \text{W/m}^2\]
\QQ{1} احسب مستوى الشدة الصوتية لصوت شدته $I=10^{-5}\ \text{W/m}^2$.
\QQ{2} بيّن أنه عند مضاعفة الشدة $I$ يزداد المستوى $L$ بحوالي $3\ \text{dB}$.
\QQ{3} آلة واحدة تُصدر صوتًا مستواه $60\ \text{dB}$. نشغّل $n$ آلة متماثلة معًا فتصبح الشدة الكلية $nI$. ما أصغر عدد $n$ من الآلات يجعل المستوى يتجاوز $85\ \text{dB}$ (عتبة الخطر على الأذن)؟

\textbf{الجزء الثاني (النمو السكاني):} يُنمذج عدد سكان مدينة بعد $t$ سنة ($t\geqslant0$) بالدالة $P(t)=12000\,e^{0{,}02t}$.
\QQ{4} احسب عدد السكان الابتدائي ثم عدد السكان بعد $10$ سنوات.
\QQ{5} بعد كم سنة يتجاوز عدد السكان $18000$ نسمة؟
\QQ{6}
\begin{sq}
\qa بيّن أنه من أجل كل $t\geqslant0$: $P(t+T)=2P(t)$ حيث $T=50\ln2$.
\qb ماذا يمثل العدد $T$؟ أعط قيمة مقربة له.
\end{sq}
\end{tamrin}

\section{اختيار من متعدد وصحيح أم خطأ}

% ===================== QCM =====================
\begin{tamrin}{اختيار من متعدد}
في كل حالة من الحالات التالية، اختر الجواب الصحيح الوحيد من بين الأجوبة الثلاثة المقترحة مع التعليل.

\QQ{1} العدد $\ln\left(e^{2}\right)-\ln\left(\sqrt{e}\right)$ يساوي:
\ABC{$2$}{$\dfrac{5}{2}$}{$\dfrac{3}{2}$}

\QQ{2} مجموعة تعريف الدالة $f:x\mapsto \ln\left(x^{2}-x\right)$ هي:
\ABC{$]0;1[$}{$]-\infty;0[\cup]1;+\infty[$}{$]1;+\infty[$}

\QQ{3} $\lim\limits_{x\to 0^{+}}x\ln x$ تساوي:
\ABC{$0$}{$-\infty$}{$+\infty$}

\QQ{4} مجموعة حلول المعادلة $\ln(x-1)+\ln(x+1)=\ln 3$ هي:
\ABC{$\{-2;2\}$}{$\{2\}$}{$\varnothing$}

\QQ{5} الدالة المشتقة للدالة $x\mapsto \ln\left(1+e^{x}\right)$ المعرفة على $\R$ هي:
\ABC{$x\mapsto \dfrac{1}{1+e^{x}}$}{$x\mapsto e^{x}\ln\left(1+e^{x}\right)$}{$x\mapsto \dfrac{e^{x}}{1+e^{x}}$}

\QQ{6} $\lim\limits_{x\to +\infty}\left(x-\ln x\right)$ تساوي:
\ABC{$+\infty$}{$0$}{$-\infty$}

\QQ{7} مجموعة حلول المتراجحة $\ln(2x-1)\leqslant 0$ هي:
\ABC{$]-\infty;1]$}{$[1;+\infty[$}{$\left]\dfrac{1}{2};1\right]$}

\QQ{8} معادلة المماس لمنحنى الدالة $\ln$ عند النقطة ذات الفاصلة $1$ هي:
\ABC{$y=x-1$}{$y=x$}{$y=x+1$}
\end{tamrin}

% ===================== صحيح أم خطأ =====================
\begin{tamrin}{صحيح أم خطأ}
أجب بـ «صحيح» أو «خطأ» مع التعليل: برّر كل عبارة صحيحة، وانقض كل عبارة خاطئة بمثال مضاد.
\QQ{1} من أجل كل عدد حقيقي $x$ غير معدوم: $\ln\left(x^{2}\right)=2\ln x$.
\QQ{2} من أجل كل عددين حقيقيين موجبين تماما $a$ و $b$: $\ln(a+b)=\ln a+\ln b$.
\QQ{3} المعادلة $\ln x=-3$ تقبل حلا وحيدا في المجال $]0;+\infty[$.
\QQ{4} من أجل كل عدد حقيقي $x$ من المجال $]0;+\infty[$: $\ln x\leqslant x-1$.
\QQ{5} $\lim\limits_{x\to+\infty}\dfrac{\ln x}{x}=1$.
\QQ{6} الدالة $x\mapsto \dfrac{\ln x}{x}$ متزايدة تماما على المجال $]0;+\infty[$.
\QQ{7} من أجل كل عدد حقيقي $x$: $\ln\left(e^{x}+1\right)>x$.
\QQ{8} إذا كان $0<a<b$ و $b\neq 1$ فإن $\dfrac{\ln a}{\ln b}<1$.
\end{tamrin}

\section{تمارين شعبة رياضيات}

% ===================== عائلة دوال =====================
\begin{tamrin}{عائلة دوال $f_m$ -- شعبة رياضيات}
$m$ وسيط حقيقي. نعتبر الدالة $f_m$ المعرفة على المجال $]0;+\infty[$ بـ: $f_m(x)=x+m\ln x$، و $(C_m)$ تمثيلها البياني في معلم متعامد ومتجانس $(O;\vec i,\vec j)$.
\QQ{1} بيّن أن جميع المنحنيات $(C_m)$ تشمل نقطة ثابتة $A$ يُطلب تعيين إحداثييها.
\QQ{2} أحسب، حسب قيم الوسيط $m$، $\lim\limits_{x\to 0^{+}}f_m(x)$ و $\lim\limits_{x\to +\infty}f_m(x)$.
\QQ{3} أدرس، حسب قيم $m$، اتجاه تغيّر الدالة $f_m$ وشكّل جدول تغيّراتها.
\QQ{4} في حالة $m<0$:
\begin{sq}
\qa بيّن أن الدالة $f_m$ تقبل قيمة حدية صغرى $k(m)=-m+m\ln(-m)$.
\qb عيّن قيم $m$ التي من أجلها يكون: من أجل كل عدد حقيقي $x$ من $]0;+\infty[$، $f_m(x)>0$.
\end{sq}
\QQ{5}
\begin{sq}
\qa اكتب معادلة للمماس $(T_m)$ للمنحنى $(C_m)$ عند النقطة $A$.
\qb عيّن قيمة $m$ التي من أجلها يكون $(T_m)$ موازيا لحامل محور الفواصل.
\end{sq}
\QQ{6} ارسم في نفس المعلم المنحنيات $(C_{-1})$ و $(C_{0})$ و $(C_{1})$.
\end{tamrin}

% ===================== مستقيم دوار =====================
\begin{tamrin}{مناقشة بيانية بمستقيم دوار -- شعبة رياضيات}
نعتبر الدالة $f$ المعرفة على المجال $]0;+\infty[$ بـ: $f(x)=x\ln x$، و $(C_f)$ تمثيلها البياني في معلم متعامد ومتجانس $(O;\vec i,\vec j)$.
\QQ{1} أحسب $\lim\limits_{x\to 0^{+}}f(x)$ و $\lim\limits_{x\to +\infty}f(x)$ و $\lim\limits_{x\to +\infty}\dfrac{f(x)}{x}$، ثم فسّر النتيجة الأخيرة بيانيا.
\QQ{2} أدرس اتجاه تغيّر الدالة $f$ ثم شكّل جدول تغيّراتها.
\QQ{3}
\begin{sq}
\qa اكتب معادلة للمماس $(T)$ للمنحنى $(C_f)$ عند النقطة $A(1;0)$.
\qb نضع من أجل كل $x>0$: $h(x)=x\ln x-x+1$. أدرس اتجاه تغيّر $h$ واستنتج إشارة $h(x)$، ثم الوضع النسبي لـ $(C_f)$ و $(T)$.
\end{sq}
\QQ{4} ارسم $(T)$ و $(C_f)$.
\QQ{5} $m$ وسيط حقيقي، و $(D_m)$ المستقيم ذو المعادلة $y=m(x-1)$.
\begin{sq}
\qa بيّن أن جميع المستقيمات $(D_m)$ تشمل النقطة $A$.
\qb ناقش بيانيا، حسب قيم الوسيط $m$، عدد حلول المعادلة $(E)$: $x\ln x-mx+m=0$ ذات المجهول $x$ من $]0;+\infty[$.
\end{sq}
\end{tamrin}

% ===================== ln x = m x =====================
\begin{tamrin}{المعادلة $\ln x=mx$ -- شعبة رياضيات}
نعتبر الدالة $g$ المعرفة على المجال $]0;+\infty[$ بـ: $g(x)=\dfrac{\ln x}{x}$، و $(C_g)$ تمثيلها البياني في معلم متعامد ومتجانس.
\QQ{1} أحسب $\lim\limits_{x\to 0^{+}}g(x)$ و $\lim\limits_{x\to +\infty}g(x)$، وفسّر النتيجتين بيانيا.
\QQ{2} أدرس اتجاه تغيّر الدالة $g$ ثم شكّل جدول تغيّراتها.
\QQ{3} حل المعادلة $g(x)=0$، واستنتج إشارة $g(x)$ حسب قيم $x$.
\QQ{4} ارسم المنحنى $(C_g)$.
\QQ{5} $m$ وسيط حقيقي.
\begin{sq}
\qa بيّن أن المعادلة $\ln x=mx$ ذات المجهول $x$ من $]0;+\infty[$ تكافئ $g(x)=m$.
\qb ناقش بيانيا، حسب قيم $m$، عدد حلول المعادلة $\ln x=mx$.
\end{sq}
\QQ{6}
\begin{sq}
\qa استنتج أنه من أجل كل عدد حقيقي $x$ من $]0;+\infty[$: $\ln x\leqslant \dfrac{x}{e}$.
\qb قارن بين العددين $e^{\pi}$ و $\pi^{e}$ (دون استعمال الحاسبة).
\end{sq}
\end{tamrin}

% ===================== وجود ووحدانية وحصر α =====================
\begin{tamrin}{وجود وحصر العدد $\alpha$ ونقطة صامدة -- شعبة رياضيات}
\textbf{I)} نعتبر الدالة $g$ المعرفة على $]0;+\infty[$ بـ: $g(x)=x+\ln x-2$.
\QQ{1} أحسب $\lim\limits_{x\to 0^{+}}g(x)$ و $\lim\limits_{x\to +\infty}g(x)$، ثم أدرس اتجاه تغيّر $g$ وشكّل جدول تغيّراتها.
\QQ{2}
\begin{sq}
\qa بيّن أن المعادلة $g(x)=0$ تقبل حلا وحيدا $\alpha$ في $]0;+\infty[$، وأن $1{,}5<\alpha<1{,}6$.
\qb تحقق أن $1{,}55<\alpha<1{,}56$.
\end{sq}
\QQ{3} استنتج إشارة $g(x)$ حسب قيم $x$.

\textbf{II)} نعتبر الدالة $f$ المعرفة على $]0;+\infty[$ بـ: $f(x)=x-2\ln x+\dfrac{1}{2}(\ln x)^{2}$.
\QQ{1} أحسب $\lim\limits_{x\to 0^{+}}f(x)$، ثم بيّن أنه من أجل كل $x>0$: $f(x)\geqslant x-2\ln x$ واستنتج $\lim\limits_{x\to +\infty}f(x)$.
\QQ{2} بيّن أنه من أجل كل عدد حقيقي $x$ من $]0;+\infty[$: $f'(x)=\dfrac{g(x)}{x}$، ثم شكّل جدول تغيّرات $f$.
\QQ{3}
\begin{sq}
\qa بيّن أن $f(\alpha)=\dfrac{\alpha^{2}}{2}+\alpha-2$.
\qb استنتج حصرا للعدد $f(\alpha)$، ثم بيّن أنه من أجل كل $x>0$: $f(x)>0$.
\end{sq}

\textbf{III)} نعتبر الدالة $\varphi$ المعرفة على $]0;+\infty[$ بـ: $\varphi(x)=2-\ln x$، و $(\Delta)$ المستقيم ذو المعادلة $y=x$.
\QQ{1} بيّن أن $\alpha$ هو الحل الوحيد للمعادلة $\varphi(x)=x$ (نقول إن $\alpha$ نقطة صامدة للدالة $\varphi$).
\QQ{2} أدرس الوضع النسبي للمنحنى $(C_\varphi)$ والمستقيم $(\Delta)$.
\QQ{3}
\begin{sq}
\qa بيّن أنه إذا كان $x\in[1{,}5;1{,}6]$ فإن $\varphi(x)\in[1{,}5;1{,}6]$.
\qb باستعمال $\varphi(\alpha)=\alpha$ والحصر $1{,}55<\alpha<1{,}56$، بيّن أن $1{,}555<\alpha<1{,}56$.
\end{sq}
\QQ{4} ارسم $(\Delta)$ و $(C_\varphi)$ في معلم متعامد ومتجانس.
\end{tamrin}

\end{document}
